% 第6章 结论

\section{结论}

本文围绕股票市场短期不稳定性的多尺度持续性及其宏观增量解释作用展开研究，以沪深300指数为样本对象，构建了以\texttt{future\_absret\_5}为主因变量、以HAR回归为基准框架的统一经验分析体系。研究的出发点，在于将经验对象从噪声更大的未来收益率方向转向未来一周平均波动幅度，从而在更具金融含义的状态变量上建立可解释且稳定的回归模型。

实证结果显示，股票市场短期不稳定性具有显著而稳定的多时间尺度持续性。过去5日和过去20日历史波动状态对未来一周平均波动幅度具有稳定解释作用，其中过去5日状态的系数最大且最为稳健，过去20日状态提供显著补充，过去60日状态的边际作用较弱。无论是在完整滚动样本还是非重叠样本中，这一结论都保持成立，说明主模型并不主要依赖于滚动窗口平滑性。就样本内外表现而言，HAR\_OLS已经形成稳定且具有较高解释力的基准模型，因此可以正式表述为"基于HAR回归的股票市场短期不稳定性基准模型"。

在扩展分析中，本文进一步将政策不确定性、汇率变化、工业品价格同比与货币增速边际变化纳入HARX框架。结果表明，宏观变量整体上未形成稳健的联合增量解释力，加入扩展变量后，训练集拟合只出现极小幅提升，样本外表现也未优于HAR基准模型。逐个进入扩展表和分样本稳定性检验显示，工业品价格同比和汇率变化在若干设定下更接近边际显著，具有一定保留价值，但总体证据仍不足以支持将宏观变量写成短期不稳定性的核心决定因素。由此可见，在当前研究对象和变量设定下，股票市场短期不稳定性的形成机制主要来自市场内部状态延续，宏观环境所起的作用更多体现为背景性扰动与补充性解释。

辅助稳健性目标\texttt{future\_logrv\_20}的结果进一步说明，波动状态指标本身具有较强持续性，且中期实现波动率的样本外拟合程度更高。不过考虑到其经济含义不如未来一周平均波动幅度直观，正式版全文仍应以\texttt{future\_absret\_5}为主目标，以\texttt{future\_logrv\_20}为稳健性辅助目标。这样的安排既兼顾统计效果，也保持了研究叙事的一致性。

综合以上结果，本文建议将正式版经验部分的主线明确为"HAR基准模型+宏观增量检验"。正文中应当以HAR\_OLS为主模型，以HARX\_OLS和HARX-lite作为扩展与补充分析，以非重叠窗口、单变量宏观扩展和分样本稳定性检验作为稳健性支撑。这样的写法更符合当前结果所支持的证据强度，也更接近正式回归论文对模型解释与结果表达的要求。
